% 对应题号 M1-71，复用已核对的本题几何图。
% 修订：启用 xeCJK 自动换行；保留原题参数与物理关系。
\documentclass[tikz,border=14pt]{standalone}
\usepackage{fontspec}
\usepackage{xeCJK}
\usepackage{amsmath,amssymb}
\usepackage{xcolor}
\usepackage{tikz}
\usetikzlibrary{arrows.meta,calc,patterns,positioning,bending}
\IfFontExistsTF{Heiti SC}{\IfFontExistsTF{Heiti SC}{\setmainfont{Heiti SC}}{\setmainfont{FandolHei-Regular}}\IfFontExistsTF{Heiti SC}{\setsansfont{Heiti SC}}{\setsansfont{TeX Gyre Heros}}}{\IfFontExistsTF{Songti SC}{\setmainfont{Songti SC}}{\setmainfont{FandolSong-Regular}}\IfFontExistsTF{Songti SC}{\setsansfont{Songti SC}}{\setsansfont{FandolSong-Regular}}}
\definecolor{bluefield}{RGB}{24,76,255}
\definecolor{redfield}{RGB}{232,49,35}
\definecolor{orangefield}{RGB}{236,132,18}
\definecolor{greenfield}{RGB}{18,132,86}
\definecolor{purplefield}{RGB}{112,72,190}
\definecolor{linegray}{RGB}{60,68,78}
\definecolor{softblue}{RGB}{238,243,255}
\definecolor{softgreen}{RGB}{238,249,244}
\definecolor{softorange}{RGB}{255,247,235}
\tikzset{
  arrow/.style={-{Latex[length=3.1mm,width=2.1mm]},line width=1.1pt},
  thinarr/.style={-{Latex[length=2.5mm,width=1.7mm]},line width=.85pt},
  axis/.style={-{Latex[length=2.5mm,width=1.7mm]},line width=.8pt,draw=linegray},
  guide/.style={draw=linegray!55,dashed,line width=.65pt},
  box/.style={draw=linegray,rounded corners=3pt,fill=white,align=center,inner sep=6pt,line width=.75pt},
  bluebox/.style={box,draw=bluefield,fill=softblue},
  greenbox/.style={box,draw=greenfield,fill=softgreen},
  orangebox/.style={box,draw=orangefield,fill=softorange}
}
\IfFontExistsTF{Heiti SC}{\setCJKmainfont[AutoFakeBold=1.5]{Heiti SC}}{\setCJKmainfont[AutoFakeBold=1.5]{FandolHei-Regular}}
\IfFontExistsTF{Heiti SC}{\setCJKsansfont{Heiti SC}}{\setCJKsansfont{FandolHei-Regular}}
\begin{document}

\begin{tikzpicture}[x=1cm,y=1cm]
\path[use as bounding box] (-10,-4.6) rectangle (10,4.6);
\fill[white] (-10,-4.6) rectangle (10,4.6);
\node[font=\Large\bfseries,text width=18.6cm,align=center] at (0,4.05) {细圆环的体密度：两个 $\delta$ 锁定子午截面点，方位角仍自由};
\begin{scope}[shift={(-4.8,-.2)}]
\draw[redfield,line width=1.5pt] (0,0) ellipse(2.6 and .9);
\draw[guide] (-3.1,0)--(3.1,0) node[right] {$x$};
\draw[guide] (0,-1.6)--(0,2.1) node[above] {$z$};
\draw[<->,bluefield] (0,0)--(2.6,0) node[midway,above] {$a$};
\node[redfield] at (0,-1.55) {$z=0$ 平面内的均匀细圆环};
\end{scope}
\begin{scope}[shift={(1.2,-.15)}]
\draw[axis] (-1.8,0)--(1.2,0) node[right] {$\rho$};
\draw[axis] (-1.5,-2)--(-1.5,2.2) node[above] {$z$};
\fill[redfield] (.75,0) circle(.1) node[above left] {$(a,0)$};
\draw[guide] (.75,-1.6)--(.75,1.6); \draw[guide] (-1.2,0)--(1.15,0);
\node[text=linegray] at (.2,-2.55) {子午截面中的支撑};
\end{scope}
\begin{scope}[shift={(6.1,-.1)}]
\node[bluebox,text width=5.6cm] at (0,1.55) {$\displaystyle \rho_{\rm v}=\frac{Q}{2\pi a}\delta(\rho-a)\delta(z)$};
\node[greenbox,text width=5.6cm] at (0,-.05) {$dV=\rho\,d\rho\,d\phi\,dz$};
\node[orangebox,text width=5.6cm] at (0,-1.7) {$\int\rho_{\rm v}dV=\frac{Q}{2\pi a}\cdot a\cdot2\pi=Q$};
\end{scope}
\end{tikzpicture}
\end{document}
